\section{The rule and the questions}\label{sec:rule}
For integers $n\ge1$ and $N\ge0$, let
\[
 X_{n,N}=\left\{a\in\Zge^n:\ \sum_{i=0}^{n-1}a_i=N\right\}.
\]
Site labels and cyclic orientation are fixed, and every index is read
modulo $n$. Put
\begin{equation}\label{eq:rule}
 q_i(a)=\min(a_i,a_{i+1}),\qquad
 F(a)_i=a_i-q_i(a)+q_{i-1}(a).
\end{equation}
All currents use the old state. Since $0\le q_i(a)\le a_i$, every updated
coordinate is nonnegative; summation telescopes, so $F$ maps $X_{n,N}$
to itself. There is no scheduler, hidden history or quotient by rotation.
The rule is an abstract finite deterministic system; no chemical
realization is asserted.

Two questions organize the note. Where does each initial state end, and
how many labelled sources produce a prescribed target in one step?
The first answer uses permanent zero separators after subtraction of the
minimum. The second uses the unique weak-ascent/strict-descent word of a
source. The temporal description includes a sharp worst clock
(Theorem~\ref{thm:clock}); the source parametrization includes every target,
not only fixed targets (Theorem~\ref{thm:atlas}). Their combination gives
an exact fixed-target product and the largest-fibre exponent
(Theorem~\ref{thm:growth}).

\paragraph{Relation to conservative rules.}
Conservation by local currents is background structure. For finite-alphabet
one-dimensional cellular automata, Boccara and Fuk\'s
\cite{boccara2002} give a local characterization of number conservation.
That characterization is not a terminal-state or all-target inverse
formula for \eqref{eq:rule}; its finite-alphabet setting also differs from
the present family over all masses $N$.

Receiver mass must be distinguished from receiver vacancy. In the
ultradiscrete Burgers rule studied by Nishinari and Takahashi
\cite{nishinari1998}, the corresponding current has the form
\[
 \min(M,a_i,L-a_{i+1}),
\]
with a capacity $L$ and an optional current bound $M$; the version without
the current bound retains $L-a_{i+1}$. Equation~\eqref{eq:rule} instead
uses $a_{i+1}$, with no fixed site-capacity parameter.
Fukuda, Segawa and Watanabe \cite{fukuda2023} derive related discrete and
ultradiscrete Burgers variants through correlated random walks; the
examined preprint formulation likewise uses a vacancy term and an
additional evolving variable. These literal comparisons do not establish
nonconjugacy under arbitrary encodings, or priority for the current.

Permanent-zero arguments, ordered piecewise-linear chambers, dyadic
integrality tests and interval products are elementary tools here.
The claims concern their fully evaluated consequences for
\eqref{eq:rule}, with scope and exclusions stated in
Section~\ref{sec:verification}.

